% Generated draft: requires mathematical and source review.
\documentclass[11pt]{article}
\usepackage[margin=25mm]{geometry}
\usepackage{fontspec}
\setmainfont{Latin Modern Roman}
\usepackage{amsmath,amssymb,mathtools,unicode-math}
\setmathfont{Latin Modern Math}
\usepackage{enumitem,needspace,xurl,hyperref}
\setlength{\parindent}{0pt}
\setlength{\parskip}{5pt}
\setlength{\emergencystretch}{3em}
\newcommand{\N}{\mathbb{N}}
\newcommand{\Z}{\mathbb{Z}}
\newcommand{\Q}{\mathbb{Q}}
\newcommand{\R}{\mathbb{R}}
\newcommand{\C}{\mathbb{C}}
\newcommand{\F}{\mathbb{F}}
\newcommand{\A}{\mathbb{A}}
\newcommand{\PP}{\mathbb{P}}
\newcommand{\E}{\mathbb{E}}
\newcommand{\eps}{\varepsilon}
\newcommand{\dd}{\,\mathrm{d}}
\begin{document}
\section*{UnsolvedMath ID: 11000270}
\textbf{Source problem code:} catalogue entry
\subsection{Definitions and mathematical statement}
Let \(K_r\) be the spine of Outer space for the rank-\(r\) free group in the source's convention, with its natural outer-automorphism action. In a standard model, Outer space parametrizes marked metric graphs and the spine is the associated forest-collapse complex. Let \(A(S)\) denote the source's complex of arc systems filling a punctured surface. The rank range, surface type, arc equivalences, and filling convention are not supplied.

\emph{Question 2.2.} Does \(K_r\) admit a compactification satisfying Rosenthal's conditions? Ask the same question for \(A(S)\). A compactification normally means a compact Hausdorff space containing the original space as a dense subspace; any required equivariance, boundary structure, and fixed-point or control properties must be specified by the actual Rosenthal conditions. Those conditions are absent from the record, so no weaker arbitrary compactification is substituted and no particular compactification is asserted to qualify.

The surrounding text says that the candidate spaces \(E\) in its preceding cases have dimension equal to the respective groups' virtual cohomological dimensions. It credits G. Mislin [68] with a cocompact classifying-space model for the mapping class group of a closed surface of the larger Teichm\"uller-space dimension. The extracted notation ``\(EG\)'' does not preserve whether a proper-action decoration is present; for groups with torsion, a classifying space for proper actions and a free contractible \(EG\) are different objects. The text then introduces a variation of Question 1.1, not supplied here. These dimension comparisons and the omitted variation are context, not a proof that either requested Rosenthal compactification exists.
\subsection{Short English statement}
Find compactifications satisfying Rosenthal's specific conditions for the Outer-space spine and for filling-arc systems on a punctured surface. The criterion is missing; the nearby minimal-dimension/proper-classifying-space discussion is retained without identifying an arbitrary compactification as sufficient.
\subsection{Sources}
\begin{itemize}
\item Local UnsolvedMath registry, permanent ID 11000270, source code catalogue entry, Question 2.2, PDF page 329. No individual proposer is supplied.
\item Farb, \emph{Problems on Mapping Class Groups and Related Topics}. Supplied URL: \url{https://www.math.uchicago.edu/~farb/papers/mcgbook.pdf}; recorded access 2026-07-29; extraction pdf-font-structure. The volume contains 329 explicitly typeset headings across chapters, with repeated labels disambiguated by page. The chapter header is ``Automorphism groups of free, surface and free abelian groups.''
\item Rosenthal's conditions are referenced without their statement or a full citation. G. Mislin [68] is credited with the cocompact mapping-class-group model; the full entry and following Question 1.1 variation are absent.
\item Supplied status open and OPEN-TRIAGE. The summary says only ``Cannot verify.'' Current primary-literature status and redistribution terms for the publicly accessible author/publisher copy remain NEEDS\_REVIEW.
\end{itemize}
\end{document}
